\documentclass[12pt]{article}

\usepackage{sbc-template}
\usepackage{graphicx,url}
\usepackage[brazil]{babel}
\usepackage[utf8]{inputenc}
\usepackage{booktabs}

\sloppy

\title{Desenvolvimento de uma Interface Grafica com PySide6 para Sistema de Locadora de Filmes}

\author{
  Aluno 1\inst{1}, Aluno 2\inst{1}, Aluno 3\inst{1}, Aluno 4\inst{1}
}

\address{
  Departamento de Computacao / Sistemas de Informacao\\
  Instituicao de Ensino Superior (IES)\\
  Disciplina: Programacao Orientada a Objetos II -- Docente: Prof. Evandro
  \email{\{aluno1, aluno2, aluno3, aluno4\}@aluno.ies.edu.br}
}

\begin{document}

\maketitle

\begin{abstract}
This paper reports the development of a graphical user interface (GUI) application designed for managing a Movie Rental Store. The software was developed in Python using the official PySide6 framework (Qt for Python), integrating all key concepts covered during the Object-Oriented Programming II course: application lifecycle and event loops, signals and slots, low-level event handling, a wide variety of widgets, layout management, toolbars, hierarchical menus, modal dialogs, message boxes, and independent additional windows. The software architecture was structured in a modular fashion, ensuring low coupling and generic code to facilitate future modifications.
\end{abstract}

\begin{resumo}
Este artigo relata o desenvolvimento de uma aplicacao de interface grafica de usuario (GUI) voltada a gestao de uma Locadora de Filmes. O software foi concebido na linguagem Python com o framework oficial PySide6 (Qt para Python), contemplando de forma integrada todos os conceitos abordados nas aulas da disciplina de Programacao Orientada a Objetos II: ciclo de vida de aplicacoes e loop de eventos, sinais e slots, manipulacao de eventos de baixo nivel, diversidade de widgets, gerenciamento de layouts, barras de ferramentas, menus hierarquicos, caixas de dialogo modais, alertas e criacao de janelas adicionais independentes. A arquitetura foi estruturada de maneira modular, com baixo acoplamento e codigo generico, viabilizando manutencoes e extensoes futuras.
\end{resumo}

\section{Introducao}

O desenvolvimento de interfaces graficas interativas e intuitivas constitui um dos pilares essenciais da Engenharia de Software e da Programacao Orientada a Objetos. No ecossistema Python, o framework Qt destaca-se como a principal tecnologia multiplataforma para construcao de aplicacoes desktop robustas.

O presente projeto tem por finalidade consolidar a aprendizagem pratica dos topicos ministrados no primeiro bloco da disciplina de Programacao Orientada a Objetos II, tendo como tema uma Locadora de Filmes. O objetivo principal consistiu em construir uma aplicacao funcional, simples e modularizada, que atenda a todos os requisitos tecnicos estipulados pelo docente:
\begin{enumerate}
    \item Utilizacao da biblioteca oficial \texttt{PySide6};
    \item Emprego do mecanismo de Sinais e Slots;
    \item Tratamento explicito de Eventos do sistema operacional/hardware;
    \item Integracao de diversos tipos de widgets;
    \item Organizacao espacial com Layouts variados;
    \item Implementacao de barras de ferramentas e menus;
    \item Aplicacao de dialogos (\texttt{QDialog}) e alertas (\texttt{QMessageBox});
    \item Criacao de pelo menos uma janela adicional independente.
\end{enumerate}

\section{Descricao Geral do Projeto e da Interface}

O sistema simula as operacoes essenciais de uma locadora de filmes, disponibilizando:
\begin{itemize}
    \item Visualizacao tabular de titulos, com informacoes de identificador (ID), titulo, genero, ano de lancamento, preco da locacao diaria, formato de midia (DVD, Blu-ray ou Digital), status de disponibilidade e avaliacao de 0 a 10;
    \item Filtragem em tempo real por titulo (busca textual) e por genero cinematografico (selecao categorica);
    \item Cadastro de novos titulos atraves de formulario modal dedicado;
    \item Exibicao de ficha tecnica completa em janela secundaria nao-bloqueante;
    \item Alternancia dinamica do status de locacao (alugar/devolver) com atualizacao instantanea da interface;
    \item Painel interativo para demonstracao e monitoramento didatico de eventos de mouse e teclado do Qt;
    \item Barra de progresso informativa refletindo a taxa de disponibilidade do acervo em tempo real.
\end{itemize}

A interface grafica e centralizada em uma janela principal do tipo \texttt{QMainWindow}, contendo uma barra de menus no topo, uma barra de ferramentas com atalhos visuais e campo de busca rapida, uma area central flexivel baseada em layout empilhado (\texttt{QStackedLayout}) e uma barra de status na parte inferior.
\section{Arquitetura do Software e Modularizacao}

Para garantir clareza, manutenibilidade e flexibilidade para customizacoes posteriores, a estrutura do projeto foi dividida em pacotes modulares:

\begin{itemize}
    \item \textbf{\texttt{main.py}}: Inicializacao da aplicacao, instanciacao de \texttt{QApplication} e execucao do loop de eventos (\texttt{app.exec()});
    \item \textbf{\texttt{dados/repositorio.py}}: Camada de armazenamento e logica de negocio dos filmes em memoria;
    \item \textbf{\texttt{componentes/dialogo\_cadastro.py}}: Subclasse de \texttt{QDialog} com botoes padrao (\texttt{QDialogButtonBox}) para insercao de filmes;
    \item \textbf{\texttt{componentes/janela\_detalhes.py}}: Subclasse de \texttt{QWidget} que opera como janela secundaria independente;
    \item \textbf{\texttt{componentes/area\_eventos.py}}: Subclasse de \texttt{QFrame} especializada na captura de eventos virtuais de baixo nivel;
    \item \textbf{\texttt{janelas/janela\_principal.py}}: Subclasse de \texttt{QMainWindow} que integra a visualizacao tabular, menus, barra de ferramentas, layouts e status bar.
\end{itemize}

\section{Documentacao dos Topicos das Aulas e Widgets Criados}

\subsection{Sinais e Slots}
O mecanismo de sinais e slots foi explorado de maneira abrangente:
\begin{itemize}
    \item Conexoes de cliques de botoes (\texttt{clicked.connect}) e acoes (\texttt{triggered.connect});
    \item Recepcao de parametros tipados em slots (\texttt{textChanged}, \texttt{currentIndexChanged}, \texttt{valueChanged});
    \item Conexao direta entre widgets (campo de busca rapida da toolbar sincronizado com o filtro da tabela);
    \item Sinais customizados com \texttt{Signal}: \texttt{statusModificado = Signal(int, bool)} para sincronizar a janela adicional com a principal, e \texttt{eventoDetectado = Signal(str)} para notificar a barra de status.
\end{itemize}

\subsection{Eventos do Sistema Operacional}
Foram sobrescritos os metodos virtuais de eventos:
\begin{itemize}
    \item \texttt{mousePressEvent(event)}, \texttt{mouseReleaseEvent(event)}, \texttt{mouseDoubleClickEvent(event)} e \texttt{mouseMoveEvent(event)}: Captura posicional e distincao de botoes do mouse;
    \item \texttt{keyPressEvent(event)}: Captura de teclas fisicas (e.g. Enter para abrir detalhes);
    \item \texttt{closeEvent(event)}: Interceptacao do fechamento da janela com dialogo de confirmacao via \texttt{QMessageBox.question}.
\end{itemize}

\subsection{Diversidade de Widgets}
Foram empregados os seguintes widgets: \texttt{QLabel}, \texttt{QPushButton}, \texttt{QLineEdit}, \texttt{QComboBox}, \texttt{QSpinBox}, \texttt{QDoubleSpinBox}, \texttt{QCheckBox}, \texttt{QRadioButton} (agrupados via \texttt{QButtonGroup}), \texttt{QSlider}, \texttt{QProgressBar}, \texttt{QTableWidget}, \texttt{QTextEdit}, \texttt{QGroupBox} e \texttt{QFrame}.

\subsection{Gerenciamento de Layouts}
A interface explora de forma aninhada:
\begin{itemize}
    \item \texttt{QVBoxLayout}: Alinhamento vertical principal;
    \item \texttt{QHBoxLayout}: Alinhamento horizontal para barras de ferramentas e filtros;
    \item \texttt{QGridLayout}: Disposicao em grade bidimensional para o formulario de cadastro;
    \item \texttt{QStackedLayout}: Alternancia em pilha entre a visao do catalogo e o painel de eventos.
\end{itemize}

\subsection{Barras de Ferramentas, Menus e Barra de Status}
A janela principal conta com \texttt{QMenuBar} com submenus e atalhos de teclado (\texttt{QKeySequence}), \texttt{QToolBar} com acoes e campo de busca embutido, \texttt{QAction} reutilizaveis e \texttt{QStatusBar} com mensagens de contexto e temporarias.

\subsection{Dialogos e Alertas (Dialogs e Alerts)}
O sistema implementa \texttt{QDialog} modal com \texttt{QDialogButtonBox} (OK/Cancelar) e multiplas modalidades de \texttt{QMessageBox}: \texttt{warning} (validacao de dados), \texttt{question} (confirmacao de exclusao e encerramento), \texttt{information} (sucesso de cadastro) e \texttt{about} (creditos).

\subsection{Janela Adicional Independente}
Conforme a Aula 04, qualquer \texttt{QWidget} sem no pai atua como uma janela independente. A classe \texttt{JanelaDetalhesFilme} foi implementada dessa forma, permitindo visualizacao assincrona/nao-bloqueante simultanea a janela principal.

\section{Registro do Uso de Inteligencia Artificial}

Em conformidade com os criterios de integridade academica estabelecidos para a avaliacao:

\begin{quote}
\textbf{Declaracao de Uso de IA:} As ferramentas de Inteligencia Artificial Generativa foram utilizadas exclusivamente para tirar duvidas pontuais de sintaxe e consultar regras acerca dos topicos das aulas para a construcao do trabalho. Toda a arquitetura do software, a logica de negocio e a redacao deste artigo foram desenvolvidas pela equipe de discentes.
\end{quote}

\section{Consideracoes Finais}

O software desenvolvido cumpre integralmente os requisitos didaticos e tecnicos estabelecidos pelo docente. A separacao clara de responsabilidades em camadas e o codigo comentado de forma pontual facilitam a realizacao de alteracoes e manutencoes posteriores pela equipe.

\begin{thebibliography}{4}
\bibitem{qtpython}
THE QT COMPANY. \textit{Qt for Python Documentation (PySide6)}. Disponivel em: \url{https://doc.qt.io/qtforpython-6/}. 2026.

\bibitem{pythonguis}
PYTHON GUIS. \textit{PySide6 Tutorials: Create GUI Applications with Python \& Qt6}. Disponivel em: \url{https://www.pythonguis.com/}. 2026.

\bibitem{evandro}
PROF. EVANDRO. \textit{Notas de Aula e Materiais da Disciplina de Programacao Orientada a Objetos II (Aulas 01 a 04)}. 2026.

\bibitem{sbc}
SOCIEDADE BRASILEIRA DE COMPUTACAO (SBC). \textit{Modelos para Publicacao de Artigos}. Disponivel em: \url{https://www.sbc.org.br/}. 2026.
\end{thebibliography}

\end{document}